%DS0



\section{ĐỀ SỐ \thesection }
\setcounter{bt}{0}
%%Câu 1.
\begin{bt} Thực hiện phép tính
	\begin{listEX}[2]
		\item $(-39)\cdot75+(-39)\cdot25$;
		\item $24:(-6)+6^2$;
		\item $[(-60):5+2\cdot(-8)]:(-2)^2+(-1)^{1024}$;
		\item $680-170:[(47-44)^2+(-2)^3\cdot2023^0]$.
	\end{listEX}
	\loigiai{
		\begin{listEX}[2]
				\item $(-39)\cdot75+(-39)\cdot25\\=(-39)\cdot(75+25)\\=(-39)\cdot100=-3900$
				\item $24:(-6)+6^2\\=-4+36\\=32$
				\item $\left[ (-60):5+2\cdot(-8)\right] :(-2)^2+(-1)^{1024}\\=(-12-16):4+1\\=-28:4+1\\=-7+1=-6$
				\item $680-170:\left[ (47-44)^2+(-2)^3\cdot2023^0\right] \\=680-170:(3^2-8\cdot1)\\=680-170:(9-8)\\=680-170:1=510$
		\end{listEX}
	}
\end{bt}


%%Câu 2.
\begin{bt} Tìm số nguyên $x$, biết:
	\begin{listEX}[2]
		\item $12\cdot x=-60$;
		\item $12\, \vdots\,  x$;
		\item $25+2\cdot(x-8)=5$;
		\item $(3x-5)-(x-2)=17$.
	\end{listEX}
	\loigiai{
		\begin{listEX}[2]
			\item $12\cdot x  = -60$\\
				\hspace*{0.72 cm}$\begin{aligned}
				x &= -60 : 12\\
				x &= -5
				\end{aligned}$
			\item 
				$12\,\vdots\,x \Rightarrow x \in \text{Ư}(12)$\\
				$x \in \{\pm1;\pm2;\pm3;\pm4;\pm6;\pm12\}$
			\item 	$25 + 2\cdot(x - 8) = 5$\\
				\hspace*{0.9 cm}$\begin{aligned}
				2\cdot(x - 8) &= -20\\
				x - 8 &= -10\\
				x &= -2
				\end{aligned}$
			\item $(3x - 5) - (x - 2)  = 17$\\
				\hspace*{2.1 cm}$\begin{aligned}
				3x - x &= 17 + 5 - 2\\
				2x &= 20\\
				x &= 10
				\end{aligned}$
		\end{listEX}
		
	}
\end{bt}



%%Câu 3.
\begin{bt} Một tòa nhà có 10 tầng (bao gồm tầng G) và 3 tầng hầm, tầng G được đánh số là tầng 0, hãy dùng phép cộng các số nguyên để diễn tả tình huống sau: Bạn Nam đang ở tầng 1, bạn sử dụng thang máy đi lên 5 tầng và sau đó đi xuống 8 tầng để lấy tài liệu. Hỏi cuối cùng, bạn Nam dừng thang máy ở tầng mấy?
	\loigiai{$1+5-8=-2$. Vậy bạn Nam đang ở tầng hầm số 2.}
\end{bt}


%%Câu 4.
\begin{bt} Một đội văn nghệ gồm 24 nam và 30 nữ được thầy giáo chia đều thành các nhóm để tập các tiết mục văn nghệ chào mừng Xuân Giáp Thìn. Hỏi:
	\begin{enumerate}
		\item Thầy giáo có thể chia nhiều nhất bao nhiêu nhóm văn nghệ?
		\item Mỗi nhóm có bao nhiêu nam, bao nhiêu nữ?
	\end{enumerate}
	\loigiai{
		\begin{enumerate}
			\item  Gọi $x$ $(x\in  \mathbb{N})$ là số nhóm nhiều nhất có thể chia. Ta có $x=\text{ƯCLN}(24,30)$\\
			Phân tích mỗi số ra thừa số nguyên tố 
				$24=2^3\cdot3$; \qquad $30=2\cdot3\cdot5$
				
			$x=\text{ƯCLN}(24,30)=2\cdot3=6$.
			
			Vậy có thể chia được nhiều nhất 6 nhóm.
			\item Mỗi nhóm có $24:6=4$ nam và $30:6=5$ nữ. 
		\end{enumerate}
	}
\end{bt}
%%Câu 5.
\begin{bt} Trường THCS S phát động chương trình quyên góp sách cũ để gửi cho trẻ em vùng sâu vùng xa. Số sách thu được đóng thành bó 40 quyển hay 50 quyển đều vừa đủ. Em hãy tính số sách quyên góp được, biết số sách quyên góp được khoảng từ 500 đến 700 quyển.
	\loigiai{Gọi $x$ $(x\in  \mathbb{N})$ là số quyển sách quyên góp được.
	Ta có $x\in \text{BC}(40,50)$ và $500<x<700$ \\
			Phân tích mỗi số ra thừa số nguyên tố 
		$40=2^3\cdot5$; \qquad $50=2\cdot5^2$.
		
		$\text{BCNN}(40,50)=2^3\cdot5^2=200 \\ \text{BC}(40,50)=B(200)=\{0;200;400;600;800;1000;...\}$.\\
		Mà $500<x<700$ nên $x=600$.\\
		Vậy có 600 quyển sách được quyên góp.}
\end{bt}

%%Câu 6.
\begin{bt} \immini[thm]{Ông Hải có một mảnh đất hình chữ nhật $ABCD$. Biết chiều rộng mảnh đất bằng 10 m và chiều dài mảnh đất bằng 20 m.
	}{
		\begin{tikzpicture}[scale=1,line width=0.8pt,line join=round,>=stealth]
			\def\a{6} \def\b{3}  \def\c{1.2} 
			\path (0,0) coordinate (B)
			(0,\b) coordinate (A)
			(\a,\b) coordinate (D)
			(\a,0) coordinate (C)
			(\c,\b) coordinate (E)
			(\a-\c,\b) coordinate (F)
			;
			\draw (A)--(B)--(C)--(D)--(A)
			(E)--(B) (F)--(C)
			(B)--(A)node[midway,left]{10 m}
			(E)--(F)node[midway,above]{12 m}
			(B)--(C)node[midway,below]{20 m}
			;
			\foreach \x/ \goc in {A/135,B/-135,C/-45,D/45,E/90,F/90}  \fill (\x) circle (0pt)        
			($(\x)+(\goc:3.5mm)$) node {$\x$};
		\end{tikzpicture}
	}
	
	\begin{enumerate}
		\item Tính diện tích mảnh đất hình chữ nhật $ABCD$.
		\item Ông Hải dành phần đất ở giữa có dạng hình thang $BEFC$ để xây nhà xưởng. Biết đáy bé hình thang là $EF=12$ m. Em hãy giúp ông Hải tính diện tích phần đất xây nhà xưởng.
		\item Phần đất còn lại ông Hải dùng để trồng rau. Biết rằng cứ mỗi mét vuông đất thu hoạch được 4 kg rau và giá bán rau là $5\,000$ đồng/kg. Tính số tiền ông Hải thu được khi bán rau.
	\end{enumerate}
	\loigiai{
		\begin{enumerate}
			\item Diện tích mảnh đất: $10\cdot20=200$ (m$^2$)
			\item Diện tích đất xây nhà xưởng: $(12+20)\cdot10:2=160$ (m$^2$)
			\item Số tiền ông Hải thu được: $(200-160)\cdot4\cdot5\,000=800\,0000$ (đồng) 
		\end{enumerate}
	}
\end{bt}


%%Câu 7.
\begin{bt} Cho tổng sau: $S=1+2+3+4+5+\ldots+2023+2024+2025$. Chứng tỏ rằng $S\, \, \vdots\, \, 3$.
	\loigiai{ $S=1+2+3+4+5+\ldots+2023+2024+2025\\=(2025+1)\cdot2025:2=2026\cdot 2025:2=2026:2\cdot 2025=1013\cdot 2025$.\\ Vì $2+0+2+5=9\, \vdots\, 3$ nên $2025\,  \, \vdots\, \,  3$. Suy ra $S\,  \, \vdots\, \,  3$.}
\end{bt}





